% ═══════════════════════════════════════════════════════════════
% SCHEME OF WORK: Spanisch Spätbeginner (Klasse 10-12)
% Dreijahresplan nach Rahmenplan MV + Lehrwerk A_tope
% ═══════════════════════════════════════════════════════════════

\documentclass[11pt, a4paper]{article}

\usepackage[utf8]{inputenc}
\usepackage[T1]{fontenc}
\usepackage[ngerman]{babel}

% --- SCHRIFTEN ---
\usepackage{helvet}
\renewcommand{\familydefault}{\sfdefault}

\usepackage{microtype}
\usepackage[margin=2cm, top=2.2cm, bottom=2cm]{geometry}
\usepackage{enumitem}
\usepackage{tabularx}
\usepackage{booktabs}
\usepackage{array}
\usepackage{colortbl}
\usepackage{multirow}
\usepackage{longtable}
\usepackage{tikz}
\usetikzlibrary{shapes.geometric, arrows.meta, positioning}

\usepackage{fancyhdr}
\usepackage{lastpage}

\usepackage{xcolor}
\usepackage{tcolorbox}
\tcbuselibrary{skins, breakable}

\usepackage{hyperref}
\hypersetup{colorlinks=true, linkcolor=spanischrot, urlcolor=spanischrot}

% ═══════════════════════════════════════════════════════════════
% FARBEN
% ═══════════════════════════════════════════════════════════════

\definecolor{spanischrot}{HTML}{C41E3A}
\definecolor{hellgrau}{HTML}{F5F5F5}
\definecolor{mittelgrau}{HTML}{666666}
\definecolor{dunkelgrau}{HTML}{333333}
\definecolor{rahmen}{HTML}{CCCCCC}
\definecolor{jahr10}{HTML}{2C5AA0}
\definecolor{jahr11}{HTML}{2A7A4B}
\definecolor{jahr12}{HTML}{B35C00}

% ═══════════════════════════════════════════════════════════════
% BOXEN
% ═══════════════════════════════════════════════════════════════

\newtcolorbox{infobox}[1][]{
    colback=hellgrau, colframe=hellgrau, boxrule=0pt, arc=0pt,
    left=10pt, right=10pt, top=8pt, bottom=8pt,
    borderline west={3pt}{0pt}{spanischrot},
    before skip=10pt, after skip=10pt, breakable
}

\newtcolorbox{jahrbox}[2][]{
    colback=white, colframe=#2, boxrule=1pt, arc=0pt,
    left=10pt, right=10pt, top=8pt, bottom=8pt,
    before skip=10pt, after skip=10pt, breakable,
    title={\textbf{#1}}, coltitle=white, colbacktitle=#2
}

% ═══════════════════════════════════════════════════════════════
% BEFEHLE
% ═══════════════════════════════════════════════════════════════

\newcommand{\es}[1]{\textit{#1}}
\newcommand{\ger}[1]{\textbf{#1}}
\newcommand{\thema}[1]{\colorbox{spanischrot!15}{\textbf{#1}}}

\setlength{\parindent}{0pt}
\setlength{\parskip}{6pt}

% Spaltentyp für zentrierte Spalten
\newcolumntype{C}[1]{>{\centering\arraybackslash}p{#1}}
\newcolumntype{L}[1]{>{\raggedright\arraybackslash}p{#1}}

% ═══════════════════════════════════════════════════════════════
% KOPF- UND FUSSZEILE
% ═══════════════════════════════════════════════════════════════

\pagestyle{fancy}
\fancyhf{}
\renewcommand{\headrulewidth}{0.5pt}
\lhead{\footnotesize\textsc{Spanisch Spätbeginner} · Klasse 10--12}
\rhead{\footnotesize Scheme of Work}
\cfoot{\footnotesize\thepage\,/\,\pageref{LastPage}}

% ═══════════════════════════════════════════════════════════════
% DOKUMENT
% ═══════════════════════════════════════════════════════════════

\begin{document}

% --- TITEL ---
\begin{center}
    {\LARGE\bfseries Scheme of Work}\\[6pt]
    {\Large Spanisch Spätbeginner (Klasse 10--12)}\\[8pt]
    {\normalsize Dreijahresplan nach Rahmenplan Mecklenburg-Vorpommern}\\[4pt]
    {\small Lehrwerk: \textit{A\_tope.com} (Cornelsen)}
\end{center}

\vspace{10pt}

% ═══════════════════════════════════════════════════════════════
% ÜBERBLICK
% ═══════════════════════════════════════════════════════════════

\section*{1. Kursstruktur und Rahmenbedingungen}

\begin{infobox}
\textbf{Besonderheit Spätbeginnerkurs:}\\[4pt]
Spanisch als neu begonnene Fremdsprache ab Klasse 10 (4 Wochenstunden). Die Schüler durchlaufen in drei Jahren den Weg von A1 bis B1(+). Ab dem zweiten Lernjahr (Klasse 11) gilt der Rahmenplan der gymnasialen Oberstufe mit den vier obligatorischen Themenfeldern.
\end{infobox}

\begin{tabularx}{\textwidth}{|L{2.5cm}|C{2cm}|C{2cm}|X|L{3cm}|}
\hline
\rowcolor{hellgrau}
\textbf{Jahrgang} & \textbf{GER-Ziel} & \textbf{Rahmenplan} & \textbf{Lehrwerk} & \textbf{Schwerpunkt} \\
\hline
\textcolor{jahr10}{\textbf{Klasse 10}} & A1 → A2 & Sek I (Kl. 7-10) & A\_tope Unidad 1-5 & Sprachliche Grundlagen \\
\hline
\textcolor{jahr11}{\textbf{Klasse 11}} & A2 → B1 & Sek II (Einf.phase) & A\_tope Unidad 6-8 + Module & Oberstufen-Themen beginnen \\
\hline
\textcolor{jahr12}{\textbf{Klasse 12}} & B1 → B1+ & Sek II (Qual.phase) & A\_tope Band 2 / Qué Pasa 3-4 & Abitur-Vorbereitung \\
\hline
\end{tabularx}

\vspace{8pt}

\subsection*{Die vier Themenfelder der Oberstufe (ab Klasse 11)}

\begin{tikzpicture}[node distance=0.3cm, every node/.style={rectangle, draw, minimum height=1cm, minimum width=5.5cm, align=center, font=\small}]
    \node[fill=spanischrot!20] (t1) {\textbf{1. El individuo y la sociedad}\\Identität, Beziehungen, Zukunft};
    \node[fill=spanischrot!20, right=of t1] (t2) {\textbf{2. La identidad nacional}\\Spanien, Lateinamerika, Kultur};
    \node[fill=spanischrot!20, below=0.4cm of t1] (t3) {\textbf{3. Política y sociedad}\\Migration, Ungleichheit};
    \node[fill=spanischrot!20, right=of t3] (t4) {\textbf{4. Desafíos modernos}\\Umwelt, Digitalisierung};
\end{tikzpicture}

\vspace{8pt}

% ═══════════════════════════════════════════════════════════════
% KLASSE 10
% ═══════════════════════════════════════════════════════════════

\newpage
\section*{2. Klasse 10 -- Erstes Lernjahr (A1 → A2)}

\begin{jahrbox}[Jahrgang 10 · ca. 160 Unterrichtsstunden · 4 WStd]{jahr10}
\textbf{Rahmenplan:} Sek I, Klasse 7-10 (adaptiert für Spätbeginner)\\
\textbf{Lehrwerk:} A\_tope.com, Unidad 1--5 (S. 8--81)\\
\textbf{Zielniveau:} A2 (Ende Klasse 10)
\end{jahrbox}

\subsection*{2.1 Sequenzübersicht Klasse 10}

\begin{longtable}{|C{0.8cm}|L{3.5cm}|L{4cm}|L{4cm}|C{1.5cm}|}
\hline
\rowcolor{jahr10!20}
\textbf{Seq} & \textbf{Thema / Lehrwerk} & \textbf{Kommunikative Ziele} & \textbf{Grammatik} & \textbf{Wochen} \\
\hline
\endfirsthead
\hline
\rowcolor{jahr10!20}
\textbf{Seq} & \textbf{Thema / Lehrwerk} & \textbf{Kommunikative Ziele} & \textbf{Grammatik} & \textbf{Wochen} \\
\hline
\endhead

\textbf{1} & \thema{¡Hola!}\newline Unidad 1 (S. 10--21) &
Sich vorstellen, begrüßen, buchstabieren, Kontaktdaten austauschen &
ser, Verben -ar/-er/-ir, Artikel, Zahlen 0-10, Verneinung &
1--4 \\
\hline

\textbf{2} & \thema{Mi gente y mi barrio}\newline Unidad 2 (S. 22--39) &
Familie beschreiben, Wohnort/Stadtviertel vorstellen, Alter angeben &
tener, estar, hay, Possessivbegleiter, Gruppenverben (e→ie, o→ue), Adjektive &
5--9 \\
\hline

\textbf{3} & \thema{¡Me gusta!}\newline Unidad 3 (S. 40--51) &
Vorlieben ausdrücken, Kleidung beschreiben, einkaufen &
gustar, estar + gerundio, Demonstrativa, indirekte Objektpronomen &
10--14 \\
\hline

& \multicolumn{4}{c|}{\cellcolor{hellgrau}\textit{Evaluación 1 -- Klassenarbeit / mündliche Prüfung}} \\
\hline

\textbf{4} & \thema{El día a día}\newline Unidad 4 (S. 52--69) &
Tagesablauf beschreiben, Uhrzeit/Datum, Schule, sich verabreden &
Imperativ, ir a + Infinitiv, reflexive Verben, saber/poder &
15--20 \\
\hline

\textbf{5} & \thema{En Madrid}\newline Unidad 5 (S. 70--81) &
Stadt beschreiben, Weg erklären, vergleichen, im Café bestellen &
Komparativ, Superlativ, direkte Objektpronomen, Ordnungszahlen &
21--26 \\
\hline

& \multicolumn{4}{c|}{\cellcolor{hellgrau}\textit{Evaluación 2 -- Klassenarbeit}} \\
\hline

\textbf{6} & \thema{Proyecto \& Repaso}\newline + Lektüre (s.u.) &
Wiederholung, Festigung, Leseförderung &
Konsolidierung aller Strukturen &
27--30 \\
\hline

\end{longtable}

\subsection*{2.2 Konkrete Lernaufgaben Klasse 10}

\begin{tabularx}{\textwidth}{|C{0.8cm}|L{4cm}|X|C{2cm}|}
\hline
\rowcolor{jahr10!20}
\textbf{Seq} & \textbf{Lernaufgabe} & \textbf{Produkt / Output} & \textbf{Sozialform} \\
\hline
\textbf{1} & \es{Mi presentación en video} -- sich selbst auf Video vorstellen (1-2 Min.) &
Video-Clip mit Name, Alter, Herkunft, Hobbys &
EA → Plenum \\
\hline
\textbf{2} & \es{El árbol genealógico} -- Stammbaum erstellen und präsentieren &
Plakat/Poster mit Familienmitgliedern + mündliche Präsentation &
EA/PA \\
\hline
\textbf{3} & \es{Desfile de moda} -- Modenschau im Klassenzimmer &
Kleidung beschreiben, Vorlieben äußern, spontanes Sprechen &
GA → Plenum \\
\hline
\textbf{4} & \es{Mi día perfecto} -- Comic oder Fotostory erstellen &
Bildgeschichte mit Sprechblasen: Tagesablauf mit Uhrzeiten &
PA/GA \\
\hline
\textbf{5} & \es{Guía de mi ciudad} -- Stadtführer erstellen &
Digitaler oder analoger Guide: Sehenswürdigkeiten, Wegbeschreibungen, Vergleiche &
GA \\
\hline
\textbf{6} & \es{Lesetagebuch} zu \textit{Gael y la red de mentiras} &
Kapitelweise Einträge: Zusammenfassung, Vokabeln, eigene Meinung &
EA \\
\hline
\end{tabularx}

\vspace{6pt}
{\footnotesize \textbf{Legende:} EA = Einzelarbeit, PA = Partnerarbeit, GA = Gruppenarbeit}

\subsection*{2.3 Ergänzende Materialien Klasse 10}

\begin{tabularx}{\textwidth}{|L{3cm}|L{4cm}|X|}
\hline
\rowcolor{hellgrau}
\textbf{Material} & \textbf{Einsatz} & \textbf{Kompetenzen} \\
\hline
\textbf{Gael y la red de mentiras}\newline (E. Rodriguez, A2) &
Lektüreprojekt Seq 6 (ca. 3-4 Wochen) &
Leseverstehen, Wortschatzerweiterung, interkulturelles Lernen (digitale Welt, Betrug online) \\
\hline
\textbf{Panorama 1}\newline (A\_tope S. 68-69) &
Integriert in Seq 4 &
Feste und Traditionen (Navidad, Nochevieja, Día de los Muertos) \\
\hline
\textbf{Kurzfilme}\newline (Youtube: cortometrajes) &
Hörsehverstehen-Übungen &
z.B. \es{``Cuerdas''} (Inklusion), \es{``El viaje de Said''} (Migration) \\
\hline
\end{tabularx}

\subsection*{2.4 Kompetenzentwicklung Klasse 10}

\begin{tabularx}{\textwidth}{|L{3cm}|X|}
\hline
\rowcolor{hellgrau}
\textbf{Kompetenz} & \textbf{Kann-Beschreibungen (A2)} \\
\hline
\textbf{Hörverstehen} & Kurze, langsam gesprochene Texte zu vertrauten Themen verstehen; Hauptaussagen erfassen \\
\hline
\textbf{Leseverstehen} & Kurze, einfache Texte (E-Mails, Dialoge, kurze Geschichten) verstehen \\
\hline
\textbf{Sprechen} & Sich vorstellen, einfache Gespräche führen, Familie und Wohnort beschreiben, Meinungen äußern \\
\hline
\textbf{Schreiben} & Kurze Texte zu vertrauten Themen verfassen (E-Mails, Beschreibungen, einfache Berichte) \\
\hline
\textbf{Sprachmittlung} & Einfache Informationen sinngemäß in die andere Sprache übertragen \\
\hline
\end{tabularx}

% ═══════════════════════════════════════════════════════════════
% KLASSE 11
% ═══════════════════════════════════════════════════════════════

\newpage
\section*{3. Klasse 11 -- Zweites Lernjahr (A2 → B1)}

\begin{jahrbox}[Jahrgang 11 · Einführungsphase · 4 WStd]{jahr11}
\textbf{Rahmenplan:} Sek II -- Gymnasiale Oberstufe (Einführungsphase)\\
\textbf{Lehrwerk:} A\_tope.com, Unidad 6--8 + fakultative Module (S. 82--178)\\
\textbf{Zielniveau:} B1 (Ende Klasse 11)\\[4pt]
\textbf{Wichtig:} Ab Klasse 11 gelten die \textbf{4 Oberstufen-Themenfelder}. Die Lehrwerksinhalte werden entsprechend zugeordnet.
\end{jahrbox}

\subsection*{3.1 Sequenzübersicht Klasse 11}

\begin{longtable}{|C{0.8cm}|L{3cm}|L{2.5cm}|L{3.5cm}|L{3cm}|C{1.2cm}|}
\hline
\rowcolor{jahr11!20}
\textbf{Seq} & \textbf{Thema} & \textbf{Themenfeld} & \textbf{Kommunikative Ziele} & \textbf{Grammatik} & \textbf{Wo.} \\
\hline
\endfirsthead
\hline
\rowcolor{jahr11!20}
\textbf{Seq} & \textbf{Thema} & \textbf{Themenfeld} & \textbf{Kommunikative Ziele} & \textbf{Grammatik} & \textbf{Wo.} \\
\hline
\endhead

\textbf{6} & \thema{Perú -- un país andino}\newline Unidad 6 &
\textcolor{spanischrot}{2. Identidad} &
Von Reisen berichten, historische Ereignisse beschreiben &
\textbf{Pretérito indefinido} (reg. + irreg.) &
1--6 \\
\hline

\textbf{7} & \thema{¿A qué te quieres dedicar?}\newline Unidad 7 &
\textcolor{spanischrot}{1. Individuo} &
Berufswünsche äußern, Fähigkeiten beschreiben, Bewerbung &
me gustaría, unpersönliche Konstr. (se, uno) &
7--11 \\
\hline

& \multicolumn{5}{c|}{\cellcolor{hellgrau}\textit{Klausur 1 (Q1) -- ersetzbar durch KLE Sprechen}} \\
\hline

\textbf{8} & \thema{Andalucía}\newline Unidad 8 &
\textcolor{spanischrot}{2. Identidad} &
Region beschreiben, früher vs. heute, Meinung äußern &
\textbf{Pretérito imperfecto}, indefinido vs. imperfecto &
12--17 \\
\hline

\textbf{9} & \thema{Vivir la diversidad}\newline Módulo 1 &
\textcolor{spanischrot}{1. Individuo} &
Wünsche/Meinungen ausdrücken &
\textbf{Subjuntivo} (Einführung) &
18--21 \\
\hline

& \multicolumn{5}{c|}{\cellcolor{hellgrau}\textit{Klausur 2 (Q2) oder KLE Sprechen}} \\
\hline

\textbf{10} & \thema{¿Te comunicas?}\newline Módulo 4 &
\textcolor{spanischrot}{4. Desafíos} &
Digitalisierung, soziale Medien, Erfahrungen berichten &
\textbf{Pretérito perfecto} &
22--25 \\
\hline

\textbf{11} & \thema{El Comercio justo}\newline Módulo 2 &
\textcolor{spanischrot}{4. Desafíos} &
Ratschläge geben, Nachhaltigkeit diskutieren &
\textbf{Condicional} &
26--29 \\
\hline

\textbf{12} & \thema{Proyecto / Film} &
\textcolor{spanischrot}{3. Política} &
Filmprojekt (s.u.), Wiederholung &
Konsolidierung &
30--32 \\
\hline

\end{longtable}

\subsection*{3.2 Konkrete Lernaufgaben Klasse 11}

\begin{tabularx}{\textwidth}{|C{0.8cm}|L{4cm}|X|C{2cm}|}
\hline
\rowcolor{jahr11!20}
\textbf{Seq} & \textbf{Lernaufgabe} & \textbf{Produkt / Output} & \textbf{Sozialform} \\
\hline
\textbf{6} & \es{Diario de viaje} -- Reisetagebuch aus Peru schreiben &
Fiktives Tagebuch (5 Einträge) mit indefinido + Bildern &
EA \\
\hline
\textbf{7} & \es{Entrevista de trabajo} -- Vorstellungsgespräch simulieren &
Rollenspiel (Arbeitgeber/Bewerber) + Feedback-Runde &
PA → Plenum \\
\hline
\textbf{8} & \es{Antes y ahora} -- Vergleich früher/heute (Region) &
Präsentation mit Bildern: Andalusien damals vs. heute &
PA/GA \\
\hline
\textbf{9} & \es{Carta a mi yo futuro} -- Brief an mein zukünftiges Ich &
Persönlicher Brief mit Wünschen + Subjuntivo &
EA \\
\hline
\textbf{10} & \es{Debate digital} -- Online-Debatte zu sozialen Medien &
Strukturierte Debatte: Pro/Contra digitale Kommunikation &
GA \\
\hline
\textbf{11} & \es{Campaña publicitaria} -- Werbekampagne für Fairtrade &
Poster/Video mit Slogans + Condicional (``Sería mejor si...'') &
GA \\
\hline
\textbf{12} & \es{Filmkritik} zu \es{``También la lluvia''} &
Schriftliche Rezension (300 Wörter) + mündliche Diskussion &
EA → Plenum \\
\hline
\end{tabularx}

\subsection*{3.3 Ergänzende Materialien Klasse 11}

\begin{tabularx}{\textwidth}{|L{3cm}|L{4cm}|X|}
\hline
\rowcolor{hellgrau}
\textbf{Material} & \textbf{Einsatz} & \textbf{Kompetenzen / Themenfeld} \\
\hline
\textbf{Kurzgeschichte}\newline \es{``Mala conexión''}\newline (A\_tope, Panorama 2) &
Seq 6: Nach Unidad 6, integriert &
Leseverstehen (A2+), Thema: Missverständnisse, Kommunikation \\
\hline
\textbf{Film:}\newline \es{``También la lluvia''}\newline (Icíar Bollaín, 2010) &
Seq 12: Filmprojekt (ca. 3-4 Wochen) &
\textbf{Themenfeld 2+3:} Kolonialismus, Wasserkrieg Bolivien, soziale Gerechtigkeit\\
Hörsehverstehen, Analyse, Diskussion \\
\hline
\textbf{Artikel:}\newline Medienberichte zu aktuellen Themen &
Kontinuierlich &
Leseverstehen, AFB II/III, Stellungnahme \\
\hline
\textbf{Panorama 3}\newline (A\_tope S. 120-125) &
Integriert in Seq 8 &
Comunidades Autónomas, regionale Identität \\
\hline
\end{tabularx}

\subsection*{3.4 Pflichtformat: Komplexe Leistungsermittlung (KLE) Sprechen}

\begin{infobox}
\textbf{Rechtliche Grundlage:} § 22 Abs. 5 APVO M-V\\[4pt]
In der Qualifikationsphase ist \textbf{mindestens eine} schriftliche Arbeit durch eine mündliche Prüfung (KLE Sprechen) zu ersetzen. Format: Paarprüfung (20 Min.), Bewertung: 60\% Sprache / 40\% Inhalt.\\[4pt]
\textbf{Empfehlung:} Durchführung in Q1 oder Q2 (Klasse 11), damit in Q3/Q4 schriftliche Übung für Abitur möglich.
\end{infobox}

% ═══════════════════════════════════════════════════════════════
% KLASSE 12
% ═══════════════════════════════════════════════════════════════

\newpage
\section*{4. Klasse 12 -- Drittes Lernjahr (B1 → B1+)}

\begin{jahrbox}[Jahrgang 12 · Qualifikationsphase · 4 WStd]{jahr12}
\textbf{Rahmenplan:} Sek II -- Gymnasiale Oberstufe (Qualifikationsphase)\\
\textbf{Lehrwerk:} A\_tope.com Band 2 \textit{(geplant)} / alternativ: Qué Pasa 3-4\\
\textbf{Zielniveau:} B1+ (Abitur)\\[4pt]
\textbf{Schwerpunkt:} Vertiefung der 4 Themenfelder, Abitur-Vorbereitung, komplexere Texte
\end{jahrbox}

\subsection*{4.1 Themenfeld-Vertiefung Klasse 12}

Da A\_tope Band 1 am Ende von Klasse 11 abgeschlossen wird, sind für Klasse 12 weitere Materialien erforderlich:

\vspace{6pt}

\textbf{Option A: A\_tope.com Band 2} (Cornelsen, ISBN 978-3-06-021338-2)

\begin{tabularx}{\textwidth}{|L{3cm}|L{4cm}|X|}
\hline
\rowcolor{hellgrau}
\textbf{Unidad} & \textbf{Thema} & \textbf{Themenfeld-Zuordnung} \\
\hline
Unidad 1 & Un mundo digital & 4. Desafíos modernos \\
\hline
Unidad 2 & Viajes por España & 2. Identidad nacional \\
\hline
Unidad 3 & Sociedad y compromiso & 3. Política y sociedad \\
\hline
Unidad 4 & El mundo laboral & 1. El individuo \\
\hline
\multicolumn{3}{|l|}{\textit{[Inhaltsverzeichnis Band 2 noch einzupflegen]}} \\
\hline
\end{tabularx}

\vspace{10pt}

\textbf{Option B: Qué Pasa 3-4} (falls Band 2 nicht genehmigt)

\begin{tabularx}{\textwidth}{|L{3cm}|X|L{3cm}|}
\hline
\rowcolor{hellgrau}
\textbf{Buch} & \textbf{Relevante Einheiten} & \textbf{Themenfeld} \\
\hline
Qué Pasa 3 & Einheiten zu Spanien, Regionen, Geschichte & 2. Identidad \\
\hline
Qué Pasa 4 & Einheiten zu Lateinamerika, aktuelle Themen & 2+3. Identidad, Política \\
\hline
\end{tabularx}

\subsection*{4.2 Projektvorschläge Klasse 12}

\begin{tabularx}{\textwidth}{|L{3cm}|L{4.5cm}|X|}
\hline
\rowcolor{hellgrau}
\textbf{Projekt} & \textbf{Beschreibung} & \textbf{Themenfeld / Kompetenzen} \\
\hline
\textbf{Literaturprojekt:}\newline \es{``La Casa''} (Paco Roca) &
Graphic Novel über Familienerinnerungen, Generationenkonflikt &
\textbf{1. El individuo:} Familie, Lebensentwürfe\\
Lesen, Analysieren, kreatives Schreiben \\
\hline
\textbf{Rechercheprojekt:}\newline \es{``Movimientos migratorios''} &
Eigenständige Recherche zu Migration Lat.Am. → Europa &
\textbf{3. Política:} Migration, soziale Realität\\
Medienkompetenz, Präsentation, Diskussion \\
\hline
\textbf{Debatte:}\newline \es{``El turismo masivo''} &
Pro-Contra-Debatte zu Massentourismus in Spanien &
\textbf{4. Desafíos:} Nachhaltigkeit, Umwelt\\
Argumentation, spontanes Sprechen \\
\hline
\end{tabularx}

\subsection*{4.3 Abitur-Vorbereitung}

\begin{itemize}[leftmargin=1.5em]
    \item \textbf{Q4-Klausur:} Unter abiturähnlichen Bedingungen (nicht ersetzbar durch KLE)
    \item \textbf{Textformate:} Sachtexte, literarische Texte, Bildanalyse, Hörverstehen
    \item \textbf{AFB-Verteilung Abitur:} 20\% AFB I, 45\% AFB II, 35\% AFB III
    \item \textbf{Stilmittel:} Metapher, Vergleich, Antithese, Hyperbel (siehe RP Sek II)
\end{itemize}

\subsection*{4.4 Konkrete Lernaufgaben Klasse 12}

\begin{tabularx}{\textwidth}{|L{4cm}|X|C{2cm}|}
\hline
\rowcolor{jahr12!20}
\textbf{Lernaufgabe} & \textbf{Produkt / Output} & \textbf{Sozialform} \\
\hline
\es{Análisis de una novela gráfica} -- \textit{La Casa} (Paco Roca) &
Schriftliche Analyse (500 Wörter): Themen, Stilmittel, eigene Stellungnahme &
EA \\
\hline
\es{Exposición oral} -- Thematische Präsentation (10 Min.) &
Strukturierte Präsentation zu einem Oberstufen-Thema + Handout &
EA → Plenum \\
\hline
\es{Ensayo argumentativo} -- Erörterung (600-800 Wörter) &
Pro-Contra-Essay zu aktuellem Thema (z.B. Tourismus, Digitalisierung) &
EA \\
\hline
\es{Simulación ONU} -- UNO-Planspiel &
Debatte zu globalem Thema (Migration, Klimawandel) als Ländervertreter &
GA \\
\hline
\es{Portafolio temático} -- Themenportfolio (Semester) &
Sammlung: Texte, Reflexionen, Vokabeln, Selbstbewertung zu einem Themenfeld &
EA \\
\hline
\end{tabularx}

% ═══════════════════════════════════════════════════════════════
% KOMPETENZPROGRESSION
% ═══════════════════════════════════════════════════════════════

\newpage
\section*{5. Kompetenzprogression über drei Jahre}

\begin{tikzpicture}[scale=0.9, transform shape]
    % Achsen
    \draw[->, thick] (0,0) -- (14,0) node[right] {Schuljahre};
    \draw[->, thick] (0,0) -- (0,7) node[above] {GER-Niveau};

    % Y-Achse Beschriftung
    \node[left] at (0,1) {A1};
    \node[left] at (0,2.5) {A2};
    \node[left] at (0,4) {B1};
    \node[left] at (0,5.5) {B1+};

    % X-Achse Beschriftung
    \node[below] at (2,0) {Kl. 10};
    \node[below] at (6,0) {Kl. 11};
    \node[below] at (10,0) {Kl. 12};
    \node[below] at (13,0) {Abitur};

    % Progression Linie
    \draw[very thick, jahr10] (0.5,0.5) -- (4,2.5);
    \draw[very thick, jahr11] (4,2.5) -- (8,4);
    \draw[very thick, jahr12] (8,4) -- (12,5);

    % Punkte
    \fill[jahr10] (0.5,0.5) circle (4pt) node[above right, font=\small] {Start};
    \fill[jahr10] (4,2.5) circle (4pt) node[above right, font=\small] {A2};
    \fill[jahr11] (8,4) circle (4pt) node[above right, font=\small] {B1};
    \fill[jahr12] (12,5) circle (4pt) node[above right, font=\small] {B1+};

    % Rahmenplan Wechsel
    \draw[dashed, gray] (4,-0.5) -- (4,6.5);
    \node[gray, font=\footnotesize, align=center] at (4,6.8) {RP Sek II\\beginnt};
\end{tikzpicture}

\vspace{12pt}

\subsection*{5.1 Sprachliche Mittel}

\begin{tabularx}{\textwidth}{|L{2.5cm}|X|X|X|}
\hline
\rowcolor{hellgrau}
\textbf{Bereich} & \textbf{Klasse 10} & \textbf{Klasse 11} & \textbf{Klasse 12} \\
\hline
\textbf{Zeiten} &
Presente &
+ Indefinido, Imperfecto, Perfecto &
+ Futuro, Condicional, Subjuntivo \\
\hline
\textbf{Syntax} &
Einfache Hauptsätze, Verneinung, Fragen &
Nebensätze (que, cuando, porque), indirekte Rede &
Komplexe Satzgefüge, Bedingungssätze \\
\hline
\textbf{Wortschatz} &
~1500 Wörter: Alltag, Familie, Schule, Stadt &
~2500 Wörter: + Berufe, Reisen, Geschichte &
~3500 Wörter: + abstrakte Begriffe, Fachvokabular \\
\hline
\end{tabularx}

\subsection*{5.2 Themenfelder-Abdeckung}

\begin{tabularx}{\textwidth}{|L{5cm}|C{2.5cm}|C{2.5cm}|C{2.5cm}|}
\hline
\rowcolor{hellgrau}
\textbf{Themenfeld} & \textbf{Kl. 10} & \textbf{Kl. 11} & \textbf{Kl. 12} \\
\hline
1. El individuo y la sociedad & (vorbereitet) & Seq 7, 9 & Vertiefung \\
\hline
2. La identidad nacional & (vorbereitet) & Seq 6, 8 & Vertiefung \\
\hline
3. Aspectos de política y sociedad & -- & Seq 12 (Film) & Schwerpunkt \\
\hline
4. Desafíos modernos & -- & Seq 10, 11 & Vertiefung \\
\hline
\end{tabularx}

% ═══════════════════════════════════════════════════════════════
% RESSOURCEN
% ═══════════════════════════════════════════════════════════════

\section*{6. Verfügbare Ressourcen}

\begin{tabularx}{\textwidth}{|L{4cm}|L{3cm}|X|}
\hline
\rowcolor{hellgrau}
\textbf{Material} & \textbf{Status} & \textbf{Einsatzbereich} \\
\hline
A\_tope.com (Band 1) & Vorhanden & Klasse 10--11 \\
\hline
A\_tope.com (Band 2) & \textcolor{jahr12}{Genehmigung ausstehend} & Klasse 12 \\
\hline
Qué Pasa 3 & Vorhanden (Klassensatz) & Alternative für Kl. 12 \\
\hline
Qué Pasa 4 & Vorhanden (Klassensatz) & Alternative für Kl. 12 \\
\hline
Gael y la red de mentiras & Vorhanden & Lektüre Kl. 10 (A2) \\
\hline
Audio-CD A\_tope & Vorhanden & Hörverstehen \\
\hline
Video-DVD A\_tope & Vorhanden & Hörsehverstehen \\
\hline
\end{tabularx}

% ═══════════════════════════════════════════════════════════════
% BINNENDIFFERENZIERUNG & INKLUSION
% ═══════════════════════════════════════════════════════════════

\section*{7. Binnendifferenzierung \& Inklusion}

\begin{infobox}
\textbf{Heterogenität im Spätbeginnerkurs:}\\[4pt]
Obwohl alle Lernenden bei Null beginnen, entstehen schnell Leistungsunterschiede durch unterschiedliche:
\begin{itemize}[leftmargin=1.5em, itemsep=2pt]
    \item Vorkenntnisse in anderen romanischen Sprachen (Französisch, Latein)
    \item Lernstrategien und Arbeitsgeschwindigkeit
    \item Motivation und Sprachbegabung
    \item Individuelle Förderbedarfe (LRS, Nachteilsausgleich)
\end{itemize}
\end{infobox}

\subsection*{7.1 Differenzierungsstrategien}

\begin{tabularx}{\textwidth}{|L{3cm}|X|X|}
\hline
\rowcolor{hellgrau}
\textbf{Strategie} & \textbf{Beschreibung} & \textbf{Beispiel} \\
\hline
\textbf{Tempo-Differenzierung} &
Zusatzaufgaben für schnelle Lernende; mehr Zeit für langsamere &
Basis: AB bearbeiten\\Erweiterung: \es{Texto creativo} verfassen \\
\hline
\textbf{Niveau-Differenzierung} &
Gleiche Aufgabe, unterschiedliche Komplexität &
$\star$ Lückentext\\$\star\star$ Sätze bilden\\$\star\star\star$ Freier Text \\
\hline
\textbf{Hilfestellung} (Scaffolding) &
Gestufte Hilfen: von Tipps bis Musterlösung &
Tipp 1: ``Welche Zeit?''\\Tipp 2: ``Unregelmäßig!''\\Tipp 3: ``\es{fue}'' \\
\hline
\textbf{Produkt-Differenzierung} &
Verschiedene Ergebnisformate möglich &
Poster ODER Podcast ODER Essay \\
\hline
\textbf{Kooperatives Lernen} &
Experten-Methode, Tandem-Lernen &
Grammatik-Experten erklären Peers \\
\hline
\end{tabularx}

\subsection*{7.2 Inklusion \& Nachteilsausgleich}

\begin{tabularx}{\textwidth}{|L{3cm}|X|}
\hline
\rowcolor{hellgrau}
\textbf{Förderbedarf} & \textbf{Mögliche Maßnahmen} \\
\hline
\textbf{LRS / Legasthenie} &
• Mehr Zeit bei schriftlichen Arbeiten (+25\%)\\
• Mündliche statt schriftliche Prüfungsanteile\\
• Größere Schrift, klares Layout\\
• Wortlisten mit Übersetzung bereitstellen \\
\hline
\textbf{Auditive Verarbeitung} &
• Texte vorab schriftlich geben\\
• Kopfhörer bei Hörverstehen\\
• Wiederholungen ermöglichen\\
• Visuelle Unterstützung (Bilder, Gestik) \\
\hline
\textbf{Konzentration / ADHS} &
• Kürzere Arbeitseinheiten\\
• Klare Struktur, Timer sichtbar\\
• Bewegungspausen einplanen\\
• Sitzplatz vorne, ablenkungsarm \\
\hline
\textbf{Hochbegabung} &
• Enrichment: komplexere Texte, authentische Materialien\\
• Compacting: Basis überspringen, Erweiterung beginnen\\
• Mentoring jüngerer Lernender \\
\hline
\end{tabularx}

\subsection*{7.3 Universal Design for Learning (UDL)}

\begin{tabularx}{\textwidth}{|L{4cm}|X|}
\hline
\rowcolor{hellgrau}
\textbf{UDL-Prinzip} & \textbf{Umsetzung im Spanischunterricht} \\
\hline
\textbf{Multiple Repräsentation}\\(Inhalte auf verschiedenen Wegen) &
• Text + Audio + Video + Bild\\
• Vokabeln: Wort + Übersetzung + Beispielsatz + Bild\\
• Grammatik: Regel + Tabelle + Farben + Beispiele \\
\hline
\textbf{Multiple Ausdrucksformen}\\(Ergebnisse auf verschiedenen Wegen) &
• Mündlich ODER schriftlich ODER multimedial\\
• Einzel ODER Partner ODER Gruppe\\
• Test ODER Portfolio ODER Präsentation \\
\hline
\textbf{Multiple Engagement}\\(Motivation auf verschiedenen Wegen) &
• Wahlmöglichkeiten bei Themen\\
• Selbsteinschätzung fördern\\
• Authentische Materialien (Songs, Filme, Influencer)\\
• Kulturelle Vielfalt abbilden (Spanien + Lat.Am.) \\
\hline
\end{tabularx}

\vspace{12pt}

\begin{center}
\rule{0.6\textwidth}{0.4pt}\\[6pt]
{\small\textit{Scheme of Work erstellt für Lehramtsprüfung 2026}}\\
{\footnotesize Grundlage: Rahmenplan Spanisch MV (Sek I + Sek II), A\_tope.com Lehrwerk}
\end{center}

\end{document}
